% BHU PYQ -- AIHC & Archaeology: Political History of North India (600 CE - 1300 CE) (2024-25 NEP)
\documentclass[11pt,a4paper]{article}
\usepackage[a4paper,top=2.5cm,bottom=2.5cm,left=2.8cm,right=2.8cm,headheight=14pt]{geometry}
\usepackage{amsmath,amssymb}
\usepackage{fontspec}
\setmainfont{Kohinoor Devanagari}
\usepackage{microtype,fancyhdr,enumitem,hyperref,lastpage,parskip}

\hypersetup{colorlinks=true,urlcolor=black,linkcolor=black,pdftitle={AIHMJ21: Political History of North India (600 CE - 1300 CE) -- BHU NEP 2024-25}}

\pagestyle{fancy}\fancyhf{}
\renewcommand{\headrulewidth}{0.4pt}\renewcommand{\footrulewidth}{0.4pt}
\fancyhead[L]{\small \textit{Banaras Hindu University}}
\fancyhead[R]{\small \textit{AIHMJ21: Political History of North India (2024-25)}}
\fancyfoot[C]{\small Page \thepage\ of \pageref{LastPage}}

\newlist{parts}{enumerate}{2}
\setlist[parts,1]{label=(\alph*),leftmargin=*,itemsep=4pt,topsep=3pt}
\setlist[parts,2]{label=(\roman*),leftmargin=*,itemsep=2pt,topsep=2pt}

\newcommand{\pts}[1]{\hfill \textit{[#1]}}

\begin{document}

\begin{center}
  {\large\bfseries BANARAS HINDU UNIVERSITY}\\[2pt]
  {\normalsize B. A. (Hons.) Semester II Examination, 2024-25 (NEP)}\\[6pt]
  \rule{0.8\linewidth}{0.5pt}\\[6pt]
  {\large\bfseries ANCIENT INDIAN HISTORY, CULTURE \& ARCHAEOLOGY}\\[4pt]
  {\large\bfseries Paper Code: AIHMJ-21 : Political History of North India (Circa 600 CE -- 1300 CE)}\\[4pt]
  {\normalsize Time: 2:30 Hours \hfill Full Marks: 70}\\[2pt]
  {\small REF NO. D-II/065}
\end{center}
\rule{\linewidth}{0.4pt}
\smallskip

\noindent \textit{Instructions: Attempt all Sections A, B and C according to instructions given. Write your Roll No.\ at the top immediately on receipt of this question paper.}

\smallskip
\noindent \textit{दिये गये निर्देशों के अनुसार खण्ड अ, ब तथा स के उत्तर दीजिये।}

\bigskip

\section*{SECTION -- A / खण्ड -- अ}
\noindent \textit{Note: Attempt any \textbf{three} questions, each in 300 words. Each question carries 10 marks.}\pts{3 $\times$ 10 = 30}

\noindent \textit{किन्हीं तीन प्रश्नों के उत्तर दीजिये, प्रत्येक उत्तर 300 शब्दों में हो। प्रत्येक प्रश्न 10 अंकों का है।}

\begin{enumerate}[label=\textbf{\arabic*.},leftmargin=*,itemsep=12pt]

  \item Write an essay on the achievements of Harshavardhana.\pts{10}
  
  हर्षवर्धन की उपलब्धियों पर एक निबन्ध लिखिए।

  \item Evaluate the political achievements of Yashovarman of Kanauj.\pts{10}
  
  कन्नौज के यशोवर्मन की राजनीतिक उपलब्धियों का मूल्यांकन कीजिए।

  \item Discuss the achievements of Mihira Bhoja I.\pts{10}
  
  मिहिर भोज प्रथम की उपलब्धियों की चर्चा कीजिए।

  \item Formulate an estimate of the career and achievements of Devapala.\pts{10}
  
  देवपाल के जीवन एवं उपलब्धियों का मूल्यांकन प्रस्तुत कीजिए।

  \item Describe the Tripartite Struggle for supremacy over Kanauj.\pts{10}
  
  कन्नौज पर आधिपत्य हेतु हुए त्रिकोणीय (त्रिपक्षीय) संघर्ष का वर्णन कीजिए।

\end{enumerate}

\bigskip

\section*{SECTION -- B / खण्ड -- ब}
\noindent \textit{Note: Attempt any \textbf{four} questions, each in 150 words. Each question carries 5 marks.}\pts{4 $\times$ 5 = 20}

\noindent \textit{किन्हीं चार प्रश्नों के उत्तर दीजिये, प्रत्येक उत्तर 150 शब्दों में हो। प्रत्येक प्रश्न 5 अंकों का है।}

\begin{enumerate}[label=\textbf{\arabic*.},leftmargin=*,start=6,itemsep=10pt]

  \item Write a note on Sasanka of Gauda.\pts{5}
  
  गौड़ राजा शशांक पर एक टिप्पणी लिखिए।

  \item Evaluate the achievements of Nagabhata II.\pts{5}
  
  नागभट्ट द्वितीय की उपलब्धियों का मूल्यांकन कीजिए।

  \item Discuss the main causes of the downfall of the Gurjara-Pratihara dynasty.\pts{5}
  
  गुर्जर-प्रतिहार वंश के पतन के मुख्य कारणों की विवेचना कीजिए।

  \item Write a short note on Chandella ruler Vidyadhara.\pts{5}
  
  चंदेल शासक विद्याधर पर एक संक्षिप्त टिप्पणी लिखिए।

  \item Write about the achievements of Paramara ruler Bhoja.\pts{5}
  
  परमार शासक भोज की उपलब्धियों के विषय में लिखिए।

  \item Discuss the invasions of Mahmud of Ghazni on North India.\pts{5}
  
  उत्तर भारत पर महमूद गजनवी के आक्रमणों की विवेचना कीजिए।

\end{enumerate}

\bigskip

\section*{SECTION -- C / खण्ड -- स}
\noindent \textit{Note: Attempt all \textbf{ten} questions of this section in a maximum of two sentences each. Each question carries 2 marks.}\pts{10 $\times$ 2 = 20}

\noindent \textit{इस खण्ड के सभी दस प्रश्नों के उत्तर अधिकतम दो वाक्यों में दीजिये। प्रत्येक प्रश्न 2 अंकों का है।}

\begin{enumerate}[label=\textbf{\arabic*.},leftmargin=*,start=12,itemsep=8pt]
  \item 
  \begin{parts}
    \item Who was the author of `Harshacharita'?\pts{2}
    
    `हर्षचरित' के रचनाकार कौन थे? (Banabhatta)

    \item Name the capital of the Pushyabhuti dynasty.\pts{2}
    
    पुष्यभूति वंश की राजधानी का नाम लिखिये। (Thanesar / Kanauj)

    \item Who was the founder of the Pala dynasty of Bengal?\pts{2}
    
    बंगाल के पाल वंश के संस्थापक कौन थे? (Gopala)

    \item Which Gurjara-Pratihara ruler adopted the title of `Adivaraha'?\pts{2}
    
    किस गुर्जर-प्रतिहार शासक ने `आदिवराह' की उपाधि धारण की थी? (Mihira Bhoja)

    \item Who was the author of `Gaudavaho'?\pts{2}
    
    `गौड़वहो' ग्रन्थ के लेखक कौन थे? (Vakpatiraja)

    \item Write the name of the capital of the Paramaras of Malwa.\pts{2}
    
    मालवा के परमारों की राजधानी का नाम लिखिये। (Dhara)

    \item Who built the Khajuraho temples?\pts{2}
    
    खजुराहो के मन्दिरों का निर्माण किसने करवाया था? (Chandellas)

    \item Name the battle fought between Prithviraj Chauhan and Muhammad Ghori in 1191 CE.\pts{2}
    
    1191 ई० में पृथ्वीराज चौहान तथा मुहम्मद गोरी के मध्य लड़े गये युद्ध का नाम लिखिये। (First Battle of Tarain)

    \item Who was the last prominent ruler of the Gahadavala dynasty of Kanauj?\pts{2}
    
    कन्नौज के गहड़वाल वंश का अन्तिम प्रमुख शासक कौन था? (Jayachandra)

    \item Which Chinese traveler visited India during the reign of King Harsha?\pts{2}
    
    राजा हर्ष के शासनकाल में किस चीनी यात्री ने भारत की यात्रा की थी? (Hiuen Tsang)
  \end{parts}
\end{enumerate}

\end{document}
